\scanpage{22}
% L1–6, one cancelled block
\deleted{$A$ est local. Montrons d'abord que les $u^i$ $(0\leq i\leq p-1)$ sont libres sur $A$. Soit en effet $\sum_0^{p-1}\lambda_i u^i=0$ une relation non-triviale entre les $u^i$, $(\lambda_i\in A)$, et soit $\lambda_\alpha$ le premier coéfficient non nul $(0\leq\alpha\leq p-1)$. Appliquant $D^\alpha$, on trouve $\unsure{\alpha!|\alpha}\,\lambda_\alpha\unsure{u|x}=0$, d'où $\lambda_\alpha=0$, absurde.}

% L7–8
Soit $u^p=a$, et soit $C'=A[\deleted{\unsure{X}}U]/(\deleted{\unsure{X}}U^p-a)$. \deleted{On \unsure{se trouve maintenant}} Soit $D'$ \inserted{une dérivation} dans $C'$ définie par $DU=1$, de sorte que $C'$ devient un espace homogène principal sous $B$ [cf ci-dessus]!
% L11–12
On a un hom \deleted{\unsure{qui}} $C'\mathrel{\unsure{\text{\textemdash}|\longrightarrow}}C$ défini $U\longrightarrow u$, et \deleted{\unsure{est}} c'est là un $B$-homomorphisme.

% L13–14
Il en résulte que \deleted{$C'\otimes C'\mathrel{\unsure{\longrightarrow|\simeq}} C\otimes C$ est} \inserted{$C'\mathrel{\text{\textemdash}} C$ est} un isomorphisme

% L15–18
[ de façon générale, si $E,E'$ sont des espaces fibrés principaux sous $G$, et \uline{si} on a un $G$-hom $E\mathrel{\text{\textemdash}}E'$, alors \deleted{\unsure{on en conclut tout de suite que|on en conclut tout de suite dans}} c'est un isomorphisme, pourvu que dans la catégorie $\unsure{\mathcal C|C}$, les \unsure{principes|diagramme} associés aux espaces à opérateurs marchent bien. Or
% L18–20, diagram between lines of the qualification
\diagram{\begin{tikzcd}[ampersand replacement=\&]
G\times E \arrow[r,dash] \arrow[d,"{\wr}"] \& G\times E' \arrow[d,"{\wr}"]\\
E\times E \arrow[r,"\sim"] \& E'\times E'
\end{tikzcd}}
% L21
\deleted{et d'autre part $E\times E\mathrel{\unsure{\text{\textemdash}|\longrightarrow}}E'\times E'$ est un}

% L22–24
C'est le cas si on se place dans la catégorie \margin{?}des anneaux finis sur $A$ \textendash\ [pourvu que $A$ soit noethérien ? \textendash\ ]
% archive sheet number: 21
